% Generated draft: requires mathematical and source review.
\documentclass[11pt]{article}
\usepackage[margin=25mm]{geometry}
\usepackage{fontspec}
\setmainfont{Latin Modern Roman}
\usepackage{amsmath,amssymb,mathtools,unicode-math}
\setmathfont{Latin Modern Math}
\usepackage{enumitem,needspace,xurl,hyperref}
\setlength{\parindent}{0pt}
\setlength{\parskip}{5pt}
\setlength{\emergencystretch}{3em}
\newcommand{\N}{\mathbb{N}}
\newcommand{\Z}{\mathbb{Z}}
\newcommand{\Q}{\mathbb{Q}}
\newcommand{\R}{\mathbb{R}}
\newcommand{\C}{\mathbb{C}}
\newcommand{\F}{\mathbb{F}}
\newcommand{\A}{\mathbb{A}}
\newcommand{\PP}{\mathbb{P}}
\newcommand{\E}{\mathbb{E}}
\newcommand{\eps}{\varepsilon}
\newcommand{\dd}{\,\mathrm{d}}
\begin{document}
\section*{UnsolvedMath ID: 30005080}
\textbf{Source problem code:} catalogue entry
\subsection{Definitions and mathematical statement}
The supplied record contains definitions and an incomplete construction, but no recoverable conjecture or complete mathematical assertion to be formalized as an open problem. The initial word ``conjecture'' has no accompanying claim.

Write \(\C^*=\C\setminus\{0\}\), and let \(X\) denote the variety carrying the \(\C^*\)-action in the fragment, with \(t\cdot x\) denoting the action. Define its fixed locus by
\[
X^{\C^*}=\{x\in X: t\cdot x=x\text{ for every }t\in\C^*\}.
\]
Interpreting the corrupted notation ``\(Y\in Y\)'' as indexing fixed point components, let \(\mathcal Y\) denote that collection. For each \(Y\in\mathcal Y\), the displayed Bia\l{}ynicki--Birula cell definitions read
\[
X^+(Y):=\left\{x\in X:\lim_{t\to0}t\cdot x\text{ exists in }X\text{ and belongs to }Y\right\},
\]
\[
X^-(Y):=\left\{x\in X:\lim_{t\to\infty}t\cdot x\text{ exists in }X\text{ and belongs to }Y\right\}.
\]
The indexing interpretation is provisional: the fragment does not define \(\mathcal Y\), specify the meaning of component, or state the geometric hypotheses and interpretation of limits needed for these definitions. In particular, no assertion that these sets are affine cells is supplied.

The fragment then describes the general question of classifying all proper geometric and semi-geometric quotients of \(G\)-invariant open subsets of a variety \(X\), where \(G\) is reductive. It reports that the case of normal proper \(X\) and \(G=\C^*\) was treated in its reference [3], using sections and semi-sections of the action and an ordered set of fixed point components. This is a supplied historical assertion, not a verified theorem statement here. The terms ``semi-geometric quotient'', ``section'', ``semi-section'', and the relevant order are not defined in the record.

The subsequent construction starts with a polarized pair \((X,L)\) carrying a nontrivial \(\C^*\)-action and a linearization, whose distinct weights on fixed point components are written
\[
a_0<a_1<\cdots<a_r.
\]
It begins the choice of a semi-section with an index \(i\in\{0,\ldots,r\}\), and the choice of a section with an index \(i\in\{0,\ldots,r-1\}\), but stops before specifying the remaining data or construction. The passage restricts attention to certain sections and semi-sections whose quotients it describes as projective standard geometric invariant theory quotients. Neither that restriction nor the quotient construction is recoverable from the supplied text. Consequently, these definitions and the construction fragment do not determine a precise open-problem statement.
\subsection{Short English statement}
The record defines attracting and repelling sets for a \(\C^*\)-action and begins a construction of quotients using linearization weights. It supplies neither a complete conjecture nor the hypotheses and construction needed to recover one. Its supplied ``open'' label is therefore not a verified mathematical status.
\subsection{Sources}
\begin{itemize}
\item Supplied local UnsolvedMath registry record, permanent ID \(30005080\), title ``Bia\l{}ynicki--Birula Cell Construction'', exact source problem code \texttt{catalogue entry}. The record belongs to the Oberwolfach Reports Open Problems collection; no proposer is supplied, and the proposed year is recorded as 2022.
\item Supplied report citation: \emph{Toric Geometry (2023)}, Oberwolfach Reports, pp.\ 861--927, DOI \texttt{10.4171/owr/2022/17}. Supplied source URL and literature-source link, labelled ``Oberwolfach Report 2022/17'': \url{https://doi.org/10.4171/owr/2022/17}. The report has not been independently consulted for this draft.
\item Literature assessment dated exactly \texttt{2026-08-22}: supplied dated curation, not independently verified. It labels the record ``open'' while assessing the extracted passage as a definition within a broader \(\C^*\)-action discussion, without a complete conjecture or verifiable problem statement. The supplied statement status is ``unrecoverable''; the supplied verification note describes the report passage as definitions and a proved construction without an open mathematical claim. These are reported source assessments, not independent verification.
\item The fragment attributes its quotient viewpoint to Mumford's Geometric Invariant Theory, refers to [2] for the Bia\l{}ynicki--Birula cells, and to [3] for the normal proper \(\C^*\)-quotient discussion. Bibliographic entries for [2] and [3], and a precise Mumford citation, are not supplied.
\end{itemize}
\end{document}
